% Omschrijven FP-LPV
% Main thing is adding C and D matrices, for output. Mapping states (6 states x and acceleration/velocity I think)
\textbf{Provide an outline of the main activities or tasks of your project and their timetable with time estimated per
activity or task. Work chronologically in months and weeks starting with day 1 of the project phase and ending
with your final presentation and defense.}
\\
\\
% Entire project
% Estimated duration
% From: 01-02-2026
% To: 27-09-2026

% Phase 2 starts
% 16-03-2026

% Deadlines:
% final presentation and defense.
% turn in paper

% Do I need to take into account my and my supervisors their vacation?

% Dual-gantry system arrives:
% also need to take into account when I can actually work on the dual-gantry system.
% Half april, misschien zelfs later eind april. Maar ben nog wel even zoet met simulatie.

% Everything should be simulation --> experiment --> real hardware setup --> validation (have to look if i make the validation experiments a subset of my experiment, I think this should be the case since I need to take it into account for experiment design.

% Onderverdelen in de volgende activiteiten aspects misschien?

% FP-LPV simulatie
% need to convert FP model to lpv state space
% need to implement FP model without LPV component/y-position dependency.
% How do I compare the both

% Continous Time First Principles model beschikbaar. Jasper heeft met ode45 (Runge–Kutta algorithm) soort van al Discrete Time. Reference en feedforward zijn ook in discrete time/hebben een bepaalde sampling rate

% LPV-LFR
% adjust code to work for mimo (basicly already done only 2 lines of code)
% need to extend Jan his code to LPV-LFR mimo
% I already have the FP-baseline (I will first extend this to be LPV-SS so that's compatible with Jan's code)
% And then I can replace the MSD variant with that. Should do that first before making it LPV-LFR etc.

% Hoe is in huidige simulatie geverifieerd


% Time planning table for Section 7
% Requires packages: booktabs, longtable, array, enumitem
% Add to preamble: \usepackage{booktabs, longtable, array, enumitem}

% ---------- Preamble ----------
% Required packages


% % Column types
% \newcolumntype{L}[1]{>{\RaggedRight\arraybackslash}p{#1}}

% % Table styling
% \renewcommand{\arraystretch}{1.16}
% \setlength{\LTpre}{0pt}
% \setlength{\LTpost}{0pt}

% % Row colors
% \definecolor{phasegray}{gray}{0.93}
% \definecolor{checkblue}{RGB}{235,242,250}
% \definecolor{deadlinepink}{RGB}{252,235,235}

% % Compact itemize for table cells
% \newlist{tightitemize}{itemize}{1}
% \setlist[tightitemize]{nosep,leftmargin=1.2em,topsep=1pt}

% % Convenience macros
% \newcommand{\periodcell}[2]{\makecell[l]{\textbf{#1}\\\footnotesize #2}}
% \newcommand{\phaserow}[1]{%
%   \rowcolor{phasegray}
%   \multicolumn{3}{L{14.4cm}}{\textbf{#1}}\\
% }
% \newcommand{\checkpointrow}[1]{%
%   \rowcolor{checkblue}
%   \multicolumn{3}{L{14.4cm}}{\textbf{Supervisor checkpoint:} #1}\\
% }
% \newcommand{\deadlinerow}[1]{%
%   \rowcolor{deadlinepink}
%   \multicolumn{3}{L{14.4cm}}{\textbf{Deadline:} #1}\\
% }

% % ---------- Table ----------
% \begin{longtable}{L{2.8cm} L{8.2cm} L{3.2cm}}
% \toprule
% \textbf{Period} & \textbf{Technical activities} & \textbf{Writing / feedback} \\
% \midrule
% \endfirsthead

% \toprule
% \textbf{Period} & \textbf{Technical activities} & \textbf{Writing / feedback} \\
% \midrule
% \endhead

% \bottomrule
% \endfoot

% \phaserow{Baseline model development}

% \periodcell{W1--W2}{Mar 16--Mar 29}
% &
% \begin{tightitemize}
%     \item Study the existing codebase and determine what interface Jan's framework expects from the baseline model
%     \item Understand Jasper's first-principles model: equations, states, inputs, outputs, parameters, and solver logic
%     \item Reproduce the current simulation setup as reference before making structural changes
% \end{tightitemize}
% &
% \begin{tightitemize}
%     \item Set up thesis structure and notation
%     \item Weekly discussion of modeling assumptions
% \end{tightitemize}
% \\
% \addlinespace[0.35em]

% \periodcell{W3--W4}{Mar 30--Apr 12}
% &
% \begin{tightitemize}
%     \item Define a discrete-time compatible baseline structure for use in the augmentation framework
%     \item Convert the first-principles model to a DT state-space form or equivalent DT update structure
%     \item Verify the DT implementation against Jasper's ODE45 simulation reference
% \end{tightitemize}
% &
% \begin{tightitemize}
%     \item Draft baseline-model implementation subsection
%     \item Record modeling choices and assumptions
% \end{tightitemize}
% \\
% \addlinespace[0.35em]

% \periodcell{W5--W6}{Apr 13--Apr 26}
% &
% \begin{tightitemize}
%     \item Refine and verify the DT baseline model
%     \item Introduce position dependence in a controlled way and formulate the LPV / quasi-LPV baseline
%     \item Compare the position-dependent baseline against frozen LTI alternatives in simulation
% \end{tightitemize}
% &
% \begin{tightitemize}
%     \item Rough draft of Aspect~1 methods
%     \item Update literature-to-method motivation
% \end{tightitemize}
% \\
% \addlinespace[0.35em]

% \phaserow{Augmentation development and hardware validation for baseline}

% \periodcell{W7--W8}{Apr 27--May 10}
% &
% \begin{tightitemize}
%     \item Study the augmentation structure in the existing framework
%     \item Replace the MSD example baseline by the DT-compatible Jasper baseline
%     \item Perform the first end-to-end simulation tests with the integrated baseline
%     \item Prepare hardware workflow, logging, and validation setup
% \end{tightitemize}
% &
% \begin{tightitemize}
%     \item Draft augmentation framework subsection
%     \item Weekly feedback on integration progress
% \end{tightitemize}
% \\
% \addlinespace[0.35em]

% \periodcell{W9--W10}{May 11--May 24}
% &
% \begin{tightitemize}
%     \item Extend the augmentation framework to the LPV-LFR setting
%     \item Verify in simulation that the augmented model improves over the baseline
%     \item Perform initial baseline checks on available hardware data or setup access
%     \item Define the first version of the operating range and experiment design
% \end{tightitemize}
% &
% \begin{tightitemize}
%     \item Rough draft of Aspect~2 methods
%     \item Short write-up of experiment-design choices
% \end{tightitemize}
% \\
% \addlinespace[0.35em]

% \checkpointrow{review baseline verification results and agree on the direction of the experimental design}
% \addlinespace[0.35em]

% \phaserow{Interpretability, experiment design, and data collection}

% \periodcell{W11--W12}{May 25--Jun 7}
% &
% \begin{tightitemize}
%     \item Implement interpretability-preserving regularisation
%     \item Compare unconstrained and constrained augmentation in simulation
%     \item Finalize the operating range, experiment categories, and held-out validation split
%     \item Prepare baseline-identification, augmentation-training, and validation experiments
% \end{tightitemize}
% &
% \begin{tightitemize}
%     \item Draft Aspect~3 methods
%     \item Draft experimental setup section
% \end{tightitemize}
% \\
% \addlinespace[0.35em]

% \checkpointrow{approve the experiment design, operating range, and validation split with ASMPT and TU/e}
% \addlinespace[0.35em]

% \periodcell{W13--W14}{Jun 8--Jun 21}
% &
% \begin{tightitemize}
%     \item Extend the regularisation approach to the relevant dual-gantry setting and document any limitations
%     \item Run baseline-related experiments on hardware
%     \item Collect augmentation-training data and held-out validation data
%     \item Familiarize with the selected black-box benchmark for comparison
% \end{tightitemize}
% &
% \begin{tightitemize}
%     \item Refine Aspect~3
%     \item Write data-collection and validation protocol
% \end{tightitemize}
% \\
% \addlinespace[0.35em]

% \phaserow{Training and validation on experimental data}

% \periodcell{W15--W16}{Jun 22--Jul 5}
% &
% \begin{tightitemize}
%     \item Inspect, clean, and structure the experimental datasets
%     \item Train the augmented model on experimental data
%     \item Train the black-box benchmark model for comparison
% \end{tightitemize}
% &
% \begin{tightitemize}
%     \item Start results chapter
%     \item Weekly feedback on first experimental outcomes
% \end{tightitemize}
% \\
% \addlinespace[0.35em]

% \checkpointrow{review the first experimental results before final validation and comparison}
% \addlinespace[0.35em]

% \periodcell{W17--W18}{Jul 6--Jul 19}
% &
% \begin{tightitemize}
%     \item Validate all models on held-out operating points and motion profiles
%     \item Compare baseline, grey-box, and black-box performance
%     \item Compare training effort, convergence behaviour, and computational cost
%     \item Use remaining time as buffer for reruns or missing experiments
% \end{tightitemize}
% &
% \begin{tightitemize}
%     \item Rough draft of Aspect~4 methods and results
%     \item Consolidate figures and tables
% \end{tightitemize}
% \\
% \addlinespace[0.35em]

% \deadlinerow{target: main experimental results and first complete comparison finalized before vacation}
% \addlinespace[0.35em]

% \phaserow{Thesis completion}

% \periodcell{W19--W20}{Jul 20--Aug 2}
% &
% Vacation
% &
% \\
% \addlinespace[0.35em]

% \deadlinerow{confirm committee composition before leaving for vacation}
% \addlinespace[0.35em]

% \periodcell{W21--W23}{Aug 3--Aug 23}
% &
% \begin{tightitemize}
%     \item Assemble all chapters into a full draft
%     \item Complete discussion, conclusion, introduction, and abstract
%     \item Check consistency between methods, results, and research questions
% \end{tightitemize}
% &
% \begin{tightitemize}
%     \item First complete thesis draft
%     \item Send draft to supervisors around Aug~21
% \end{tightitemize}
% \\
% \addlinespace[0.35em]

% \periodcell{W24--W25}{Aug 24--Sep 6}
% &
% \begin{tightitemize}
%     \item Incorporate supervisor feedback from weekly meetings
%     \item Clean up implementation code, validation scripts, and workflow documentation
%     \item Finalize appendices and reproducibility details
% \end{tightitemize}
% &
% \begin{tightitemize}
%     \item Revision round based on supervisor comments
%     \item Improve discussion and conclusion
% \end{tightitemize}
% \\
% \addlinespace[0.35em]

% \periodcell{W26}{Sep 7--Sep 13}
% &
% \begin{tightitemize}
%     \item Final thesis revisions
%     \item Final checks on formatting, references, and confidentiality-sensitive content
% \end{tightitemize}
% &
% \begin{tightitemize}
%     \item Finalize submission version
% \end{tightitemize}
% \\
% \addlinespace[0.35em]

% \deadlinerow{complete the confidentiality procedure by approximately Sep~11}
% \addlinespace[0.35em]

% \periodcell{W27--W28}{Sep 14--Sep 27}
% &
% \begin{tightitemize}
%     \item Share the thesis with the committee
%     \item Prepare the final presentation
%     \item Rehearse the defense and refine the presentation story line
% \end{tightitemize}
% &
% \begin{tightitemize}
%     \item Final polishing after committee sharing
%     \item Prepare defense slides
% \end{tightitemize}
% \\
% \addlinespace[0.35em]

% \deadlinerow{share the thesis with the committee by Sep~27}
% \addlinespace[0.35em]

% \periodcell{After Sep 27}{TBD}
% &
% \begin{tightitemize}
%     \item Final presentation and defense
%     \item Send the thesis and signed TU/e code of scientific conduct according to the graduation procedure
% \end{tightitemize}
% &
% \\

% \end{longtable}



% Time planning table for Section 7
% Requires packages: booktabs, longtable, array, enumitem, makecell, xcolor, ragged2e
% Preamble definitions:

% Column types
\newcolumntype{L}[1]{>{\RaggedRight\arraybackslash}p{#1}}

% Table styling
\renewcommand{\arraystretch}{1.16}
\setlength{\LTpre}{0pt}
\setlength{\LTpost}{0pt}

% Row colors
\definecolor{phasegray}{gray}{0.93}
\definecolor{checkblue}{RGB}{235,242,250}
\definecolor{deadlinepink}{RGB}{252,235,235}

% Compact itemize for table cells
\newlist{tightitemize}{itemize}{1}
\setlist[tightitemize]{nosep,leftmargin=1.2em,topsep=1pt}

% Convenience macros
\newcommand{\periodcell}[2]{\makecell[l]{\textbf{#1}\\\footnotesize #2}}
\newcommand{\phaserow}[1]{%
  \rowcolor{phasegray}
  \multicolumn{3}{L{14.4cm}}{\textbf{#1}}\\
}
\newcommand{\checkpointrow}[1]{%
  \rowcolor{checkblue}
  \multicolumn{3}{L{14.4cm}}{\textbf{Supervisor checkpoint:} #1}\\
}
\newcommand{\deadlinerow}[1]{%
  \rowcolor{deadlinepink}
  \multicolumn{3}{L{14.4cm}}{\textbf{Deadline:} #1}\\
}

% ---------- Table ----------
\begin{longtable}{L{2.8cm} L{8.2cm} L{3.2cm}}
\toprule
\textbf{Period} & \textbf{Technical activities} & \textbf{Writing / feedback} \\
\midrule
\endfirsthead

\toprule
\textbf{Period} & \textbf{Technical activities} & \textbf{Writing / feedback} \\
\midrule
\endhead

\bottomrule
\endfoot

%% ==== PHASE 1: BASELINE (3 weeks) ====
\phaserow{Baseline model development}

\periodcell{W1--W2}{Mar 16--Mar 29}
&
\begin{tightitemize}
    \item Study the existing augmentation structure to determine what structure the baseline model needs to satisfy
    \item Understand Jasper's first-principles model
    \item Discretize the first-principles model to a DT state-space form (discretization method to be motivated)
    \item Verify the DT model outputs against Jasper's ODE45 simulation reference
\end{tightitemize}
&
\begin{tightitemize}
    \item Set up thesis structure and notation
    \item Record modeling choices and assumptions
\end{tightitemize}
\\
\addlinespace[0.35em]

\periodcell{W3}{Mar 30--Apr 5}
&
\begin{tightitemize}
    \item Verify that the quasi-LPV interpretation holds
    \item Verify inertia matrix invertibility across the operating range 
    \item Create frozen LTI alternatives by fixing $Y$ and compare against the position-dependent baseline in simulation
\end{tightitemize}
&
\begin{tightitemize}
    \item Rough draft of Aspect~1 methods
\end{tightitemize}
\\
\addlinespace[0.35em]

\checkpointrow{review that the DT baseline model is verified and discuss the LPV vs LTI comparison results}
\addlinespace[0.35em]

%% ==== PHASE 2: AUGMENTATION + HARDWARE (5 weeks) ====
\phaserow{Augmentation development and hardware validation for baseline}

\periodcell{W4--W5}{Apr 6--Apr 19}
&
\begin{tightitemize}
    \item Integrate the first-principles baseline into the augmentation framework
    \item Train the first augmentation on simulated data (fixed $Y$, no scheduling)
    \item Verify the augmented model improves over the baseline on simulated settling trajectories
\end{tightitemize}
&
\begin{tightitemize}
    \item Draft augmentation framework subsection
\end{tightitemize}
\\
\addlinespace[0.35em]

\periodcell{W6--W8}{Apr 20--May 10}
&
\begin{tightitemize}
    \item Extend the augmentation to the LPV-LFR setting (scheduling on $Y$)
    \item Commission dual-gantry hardware (expected mid-to-late April)
    \item Validate the baseline on hardware and identify where it fails
    \item Begin defining the operating range and experiment design, informed by the validation results
\end{tightitemize}
&
\begin{tightitemize}
    \item Rough draft of Aspect~2 methods
    \item Initial experiment-design write-up
\end{tightitemize}
\\
\addlinespace[0.35em]

\checkpointrow{review augmentation simulation results and baseline hardware validation; agree on experiment design direction}
\addlinespace[0.35em]

%% ==== PHASE 3: INTERPRETABILITY + EXPERIMENTS (5 weeks) ====
\phaserow{Interpretability and experimental design}

\periodcell{W9--W10}{May 11--May 24}
&
\begin{tightitemize}
    \item Implement orthogonal projection-based regularisation
    \item Compare constrained and unconstrained augmentation in simulation
    \item Finalize experiment design and validation split with ASMPT supervisors
\end{tightitemize}
&
\begin{tightitemize}
    \item Draft Aspect~3 methods
    \item Draft experimental setup section
\end{tightitemize}
\\
\addlinespace[0.35em]

\checkpointrow{approve the experiment design, operating range, and validation split with ASMPT and TU/e}
\addlinespace[0.35em]

\periodcell{W11--W13}{May 25--Jun 14}
&
\begin{tightitemize}
    \item Extend the regularisation to the dual-gantry setting (MIMO, position-dependent, LFR); document limitations where needed
    \item Collect augmentation-training data and held-out validation data
\end{tightitemize}
&
\begin{tightitemize}
    \item Refine Aspect~3 write-up
    \item Write data-collection protocol
\end{tightitemize}
\\
\addlinespace[0.35em]

%% ==== PHASE 4: TRAINING + VALIDATION (5 weeks) ====
\phaserow{Training and validation on experimental data}

\periodcell{W14--W15}{Jun 15--Jun 28}
&
\begin{tightitemize}
    \item Prepare experimental data for training
    \item Train the augmented model on experimental data with orthogonality regularisation
    \item Set up and train the SUBNET as black-box comparison model
\end{tightitemize}
&
\begin{tightitemize}
    \item Start results chapter
\end{tightitemize}
\\
\addlinespace[0.35em]

\checkpointrow{review first experimental results before final validation and comparison}
\addlinespace[0.35em]

\periodcell{W16--W17}{Jun 29--Jul 12}
&
\begin{tightitemize}
    \item Validate all models on held-out operating points and motion profiles; report BFR
    \item Compare training efficiency and computational cost: grey-box vs black-box
\end{tightitemize}
&
\begin{tightitemize}
    \item Draft Aspect~4 methods and results
    \item Consolidate figures and tables
\end{tightitemize}
\\
\addlinespace[0.35em]

\periodcell{W18}{Jul 13--Jul 19}
&
\begin{tightitemize}
    \item Buffer for reruns, missing experiments, or refinement
\end{tightitemize}
&
\begin{tightitemize}
    \item Finalize results write-up
\end{tightitemize}
\\
\addlinespace[0.35em]

\checkpointrow{confirm experimental results are complete before vacation}
\addlinespace[0.35em]

%% ==== PHASE 5: THESIS COMPLETION (7 working weeks + 2 vacation) ====
\phaserow{Thesis completion}

\periodcell{W19--W20}{Jul 20--Aug 2}
&
Vacation
&
\\
\addlinespace[0.35em]

\deadlinerow{confirm committee composition with supervisor before leaving for vacation}
\addlinespace[0.35em]

\periodcell{W21--W23}{Aug 3--Aug 23}
&
\begin{tightitemize}
    \item Assemble all chapters into a full thesis draft
    \item Write discussion, conclusion, introduction, and abstract
    \item Check consistency between methods, results, and research questions
\end{tightitemize}
&
\begin{tightitemize}
    \item First complete thesis draft
    \item Send draft to supervisors around Aug~21
\end{tightitemize}
\\
\addlinespace[0.35em]

\periodcell{W24--W25}{Aug 24--Sep 6}
&
\begin{tightitemize}
    \item Incorporate supervisor feedback from weekly meetings
    \item Clean up implementation code, validation scripts, and workflow documentation for ASMPT
    \item Finalize appendices and reproducibility details
\end{tightitemize}
&
\begin{tightitemize}
    \item Revision round based on supervisor comments
    \item Improve discussion and conclusion
\end{tightitemize}
\\
\addlinespace[0.35em]

\periodcell{W26}{Sep 7--Sep 13}
&
\begin{tightitemize}
    \item Final thesis revisions
    \item Final checks on formatting, references, and confidentiality-sensitive content
\end{tightitemize}
&
\begin{tightitemize}
    \item Finalize submission version
\end{tightitemize}
\\
\addlinespace[0.35em]

\deadlinerow{complete the confidentiality procedure by approximately Sep~11}
\addlinespace[0.35em]

\periodcell{W27--W28}{Sep 14--Sep 27}
&
\begin{tightitemize}
    \item Share the thesis with the committee
    \item Prepare the final presentation
    \item Rehearse the defense and refine the presentation
\end{tightitemize}
&
\begin{tightitemize}
    \item Final polishing after committee sharing
    \item Prepare defense slides
\end{tightitemize}
\\
\addlinespace[0.35em]

\deadlinerow{share the thesis with the committee by Sep~27}
\addlinespace[0.35em]

\periodcell{After Sep 27}{TBD}
&
\begin{tightitemize}
    \item Final presentation and defense
    \item Send the thesis and signed TU/e code of scientific conduct according to the graduation procedure
\end{tightitemize}
&
\\

\end{longtable}